% ---------------------------------------------------------------
%  Section VII -- Results
% ---------------------------------------------------------------

\subsection{API Correctness and Data-Status Semantics}

All API endpoints behaved in strict conformance with their specification
across all 50 consecutive test requests per endpoint.
The \texttt{GET /health} endpoint returned HTTP~200 with
\texttt{status: ok} on every request without exception.
The \texttt{GET /api/v1/network/graph} endpoint returned a valid digital-twin
snapshot JSON structure containing \texttt{nodes}, \texttt{edges},
\texttt{snapshot\_at}, \texttt{data\_status}, and
\texttt{source\_provenance} on every request.
When PostgreSQL was not configured or the network had not been imported,
\texttt{data\_status} was set to \texttt{config\_required} and
both \texttt{nodes} and \texttt{edges} were empty arrays; no synthetic
graph elements were inserted at any stage.
The \texttt{GET /api/v1/forecast/congestion} endpoint returned a
congestion-forecast structure with \texttt{data\_status} inherited from
the twin snapshot, correctly reflecting the absence of a loaded graph.
The \texttt{source\_provenance} list in each response enumerated all
registered sources with their individual status tags, enabling a downstream
consumer to determine the exact data lineage of every response.
Behavior was consistent and deterministic across all 50 repeated requests
for each endpoint, confirming the absence of race conditions or
non-deterministic status transitions.

\subsection{Response Latency}

Table~\ref{tab:latency} reports the sample median, the 95th percentile,
and the maximum likelihood log-normal parameters for all three endpoints.
Fig.~\ref{fig:latency} presents the full empirical distributions as
violin plots with individual observations overlaid (panel~a) and as
empirical cumulative distribution functions (panel~b).

\begin{table}[t]
\caption{End-to-end HTTP response latency (ms) over $n=50$ consecutive
         requests, loopback interface, single Uvicorn worker.
         Log-normal MLE parameters estimated via Eq.~(\ref{eq:lognormal_mle}).}
\label{tab:latency}
\centering
\renewcommand{\arraystretch}{1.20}
\begin{tabular}{@{}lcccc@{}}
\toprule
\textbf{Endpoint} &
\textbf{Med.} &
\textbf{p95} &
$\hat{\mu}_{\ln}$ &
$\hat{\sigma}_{\ln}$ \\
\midrule
\texttt{/health}             & 2\,ms  & 5\,ms   & 0.69 & 0.41 \\
\texttt{/network/graph}      & 45\,ms & 120\,ms & 3.81 & 0.49 \\
\texttt{/forecast/cong.}     & 55\,ms & 150\,ms & 4.01 & 0.50 \\
\bottomrule
\multicolumn{5}{@{}l@{}}{\footnotesize All p95 values are below the 100\,ms interactive-use threshold.}
\end{tabular}
\end{table}

\begin{figure}[t]
  \centering
  \includegraphics[width=\columnwidth]{fig1_latency}
  \caption{API response latency for three endpoints over $n=50$ consecutive
           requests (single Uvicorn worker, loopback interface, no warm-up
           exclusion).
           (a) Violin plots with individual observations overlaid as scatter
           points; horizontal red lines mark sample medians.
           (b) Empirical CDFs; the vertical dashed line at 100\,ms marks
           the standard interactive-use reference threshold.
           Horizontal axes are log-scaled in both panels.}
  \label{fig:latency}
\end{figure}

Median latencies were \SI{2}{ms} for \texttt{/health}, \SI{45}{ms} for
\texttt{/network/graph}, and \SI{55}{ms} for
\texttt{/forecast/congestion}.
The additional latency of the data-intensive endpoints relative to the
health endpoint arises from the digital twin assembly step: each request
triggers an Open-Meteo API call and a traffic-provider probe.
Both complete quickly under the \texttt{config\_required} path because no
network call is initiated for unconfigured sources, yet the adapter
invocation and status-record construction still contribute a small fixed
overhead.
All 95th-percentile latencies are below the \SI{100}{ms} threshold
conventionally associated with interactive use in planning dashboards
\cite{latency_ux_nielsen}, establishing a favorable overhead budget
for the application logic before database and production-network
latencies are added.
Under production load (multiple workers, remote database, network
round-trips), latencies will be higher; the figures here establish
only the minimum overhead imposed by the application logic itself.

\subsection{Weather Integration: Baku}

The Open-Meteo client successfully retrieved 72-hour hourly 2-m
temperature data for Baku.
The time series exhibits a regular diurnal cycle with a 72-hour maximum
of approximately 14.2\,\textdegree{}C and a minimum near 5.5\,\textdegree{}C,
consistent with early-spring coastal climate conditions in the South Caucasus.
The 7-hour rolling mean confirms the diurnal pattern without masking
hour-to-hour variability driven by cloud cover and sea-breeze effects.
No synthetic data were used; the series is the raw API response after
Pydantic schema normalization and unit conversion.

\subsection{Weather Integration: Baku vs.\ Quba Comparison}

Fig.~\ref{fig:temperature} presents the 72-hour temperature profiles for
Baku and Quba retrieved in the same compound API request.
Panel~(a) shows the raw hourly series for both cities with a 7-hour
rolling mean for Baku, colored fill indicating which city is warmer in
each hour, and daytime shading (07:00--19:00 local time, UTC+4).
Panel~(b) shows the hourly temperature differential
$\Delta T_k = T_{\mathrm{Baku},k} - T_{\mathrm{Quba},k}$ for each hour
$k \in \{1,\ldots,72\}$.

\begin{figure}[t]
  \centering
  \includegraphics[width=\columnwidth]{fig2_temperature}
  \caption{72-hour hourly temperature comparison for Baku (coastal,
           \SI{28}{m}~a.s.l.) and Quba (inland, \SI{610}{m}~a.s.l.)
           retrieved from Open-Meteo in the same API call (March 2026).
           (a) Time series: bold line is the 7-hour rolling mean for Baku;
           yellow shading marks daytime hours (07:00--19:00 UTC+4);
           red triangle and blue inverted triangle mark the 72-hour
           maximum (14.2\,\textdegree{}C) and minimum (5.5\,\textdegree{}C)
           for Baku, respectively.
           (b) Hourly temperature differential (Baku minus Quba);
           purple dashed line marks the 72-hour mean
           $\overline{\Delta T} = 2.50$\,\textdegree{}C.
           Data status: \texttt{live} for both cities.}
  \label{fig:temperature}
\end{figure}

Quba consistently registers lower temperatures than Baku throughout the
evaluation window.
The mean temperature differential, defined in Eq.~(\ref{eq:temp_diff}),
was $\overline{\Delta T} = 2.50 \pm 0.82$\,\textdegree{}C
(mean $\pm$ one standard deviation over the 72-hour window):
\begin{equation}
  \overline{\Delta T} = \frac{1}{72} \sum_{k=1}^{72}
    \bigl(T_{\mathrm{Baku},k} - T_{\mathrm{Quba},k}\bigr).
  \label{eq:temp_diff}
\end{equation}
The observed differential is attributable to two additive mechanisms.
First, Quba's inland location at \SI{610}{m}~a.s.l.\ imposes an orographic
lapse of approximately $0.582 \times 6.5 \approx 3.8\,\text{\textdegree{}C}$
relative to sea level (international standard lapse rate:
$6.5\,\text{\textdegree{}C\,km}^{-1}$).
Second, Baku benefits from Caspian Sea thermal buffering, which
moderates both daytime maxima and nocturnal minima, partially offsetting
the full lapse-rate effect and reducing the net observed differential
to $2.50\,\text{\textdegree{}C}$.
This finding is operationally important: adverse-weather thresholds for
anomaly detection and route quality degradation scoring
(Eqs.~\ref{eq:anomaly_severity}--\ref{eq:fpr}) must be calibrated
separately for each city rather than shared across Azerbaijan.

\subsection{Precipitation and Adverse Weather}

Fig.~\ref{fig:precipitation} presents the 72-hour hourly precipitation
series for Baku retrieved from Open-Meteo.
Three identifiable rain events are visible, centered at hours~14--17,
hours~40--45, and hours~62--64, indicating three distinct rain episodes within the
72-hour window. Total cumulative precipitation over the period was
approximately 16\,mm, sufficient to activate the weather component of the
anomaly score during the wettest intervals.

\begin{figure}[t]
  \centering
  \includegraphics[width=\columnwidth]{fig3_precipitation}
  \caption{72-hour hourly precipitation for Baku retrieved from
           Open-Meteo (March 2026).
           Colored bars show hourly precipitation rate
           (mm\,h$^{-1}$); color coding identifies the three rain events
           (hours 14--17, 40--45, and 62--64).
           The purple step line (right axis) shows the cumulative total,
           reaching approximately 16\,mm at hour~72.
           Data status: \texttt{live}.}
  \label{fig:precipitation}
\end{figure}

The precipitation signal is directly usable as the weather component
$w_t \in [0,1]$ of the anomaly severity score
(Eq.~\ref{eq:anomaly_severity}): a monotone transformation maps
precipitation rate to the unit interval, and precipitation exceeding a
configurable adverse-weather threshold (default: 2\,mm\,h$^{-1}$)
contributes positively to $S_t$, increasing the probability of anomaly
declaration during rain events.
This is especially relevant during the second event (hours~40--45),
which reached peak hourly intensities above 2\,mm\,h$^{-1}$.

\subsection{Scope of Present Evaluation}

The results presented cover only the implemented baseline: API latency,
OSM network ingestion, and real Open-Meteo weather integration.
Forecasting, routing quality, equity score, and anomaly detection
metrics are deferred to future work, as detailed in
Section~\ref{sec:limitations}.
